\documentclass[11pt,a4paper]{article}

% --- Paquets pour la rigueur scientifique ---
\usepackage[utf8]{inputenc}
\usepackage[T1]{fontenc}
\usepackage[english]{babel}
\usepackage{amsmath, amssymb, amsfonts, amsthm}
\usepackage{geometry}
\geometry{margin=2.5cm}
\usepackage{hyperref}
\usepackage{booktabs}
\usepackage{graphicx}
\usepackage{enumitem}
\setcounter{tocdepth}{1}

% --- Métadonnées ---

\title{The Emergence of Physical Reality\\
within the Framework of the Projective Dynamic Logo (PDL)}

\author{Cédric Laubscher\\
\small Independent Researcher}

\date{February 2, 2026}

\begin{document}
\maketitle

\section*{Teaser}

The majority of approaches to the “foundations of physics” begin by positing a pre-existing space–time background, a specific dynamical law, or a probabilistic framework, and subsequently attempt to account for why the physical world exhibits the particular spectrum of particles and parameter values that we empirically observe. In contrast, the Projective Dynamic Logo (PDL) adopts an inverse strategy: it proceeds from an almost minimal starting point—namely, four purely relational axioms formulated on signed graphs—and then investigates which mathematical and physical structures are logically compelled to emerge, under the assumption that anything is to exist in a temporally persistent manner at all.

From this point of departure, PDL undertakes three distinct operations that, in combination, differentiate it from the majority of prior foundational frameworks.

\begin{enumerate}
    \item \textbf{It derives an elementary ``particle'' through formal logical analysis rather than through conjecture.}

    PDL represents reality as a network of minimal logical closures, formally described as finite signed graphs whose edges encode binary agreement/disagreement relations. The evolution of these configurations is governed by a coherence–optimisation principle, rather than by externally prescribed dynamical laws. Under four axioms—minimal binary pulsation, triangular coherence, minimal completeness, and logical optimisation—it can be shown that the first admissible elementary stationary structure is a complete signed graph on four vertices and six edges, referred to as the \emph{(4,6) structure}. There exists a particular assignment of signs on the edges such that all triangular subgraphs satisfy the coherence condition, and such that a global pulsation $s \leftrightarrow -s$ leaves this coherence invariant. This structure is not \emph{a priori} postulated to “be” an electron; rather, it is identified as the minimal self-sustaining regime of existence admitted by the axioms. The electron at rest is then interpreted as the first physical realisation of this logically stationary regime, with its Compton pulsation corresponding to the internal 2–cycle, and $\hbar$ interpreted as the minimal quantum of action associated with a complete coherence cycle.

   \item \textbf{It establishes a precise combinatorial framework for modeling the proton, in place of a merely qualitative or heuristic description.}

    PDL does not treat the proton as an opaque effective degree of freedom, nor merely as an abstract three–quark bound state characterized solely by its global quantum numbers. Instead, it posits an explicit two–tier structural architecture: three localized stationary domains (“valence cores”), each represented as a superposition of logical tetrads (with $n_u = 24$ and $n_d = 28$), embedded within a dense relational sea containing $\mathcal{O}(10^4)$ dynamical relations, and terminating at a finite active boundary that mediates the coupling to an electron–type $(4,6)$ structure. 

    This construction is not introduced in an ad hoc manner: the numerical values (24, 28, 930, 10 087, 11 017) are derived from the same optimisation functional that selects the $(4,6)$ block, subject to explicit combinatorial constraints (multiplicity in units of 4, near-maximal triangular coherence, net charge $+1$, and a prescribed stationary/dynamical partition). The resulting proton configuration displays qualitative analogies with quantum chromodynamics (QCD), particularly in that it yields a decomposition of the mass into approximately $9\%$ stationary and $91\%$ dynamical components, while remaining entirely formulated within a discrete, graph-theoretic framework.

    \item \textbf{It associates multiple dimensionless constants with a single relational substratum, characterized by nontrivial numerical structure.}

    Within this architectural framework, conventionally recognized constants assume a unified structural significance, rather than functioning as independent empirically determined input parameters:
    \begin{itemize}[label=-]
        \item The reduced Planck constant $\hbar$ is interpreted as quantifying the minimal action associated with a single coherence cycle of the elementary $(4,6)$ dynamical regime.
        \item The fine–structure constant $\alpha$ is reinterpreted as a surface–to–volume coherence fraction, namely the proportion of the proton’s total coherence “budget” that is effectively projected onto its active surface and thus available for electron binding. This quantity is expressible in terms of the proton–electron mass ratio $\mu$, the valence content $r_{\text{valence}} = 930$, and the golden ratio $\varphi$, through an effective active surface $R^{(\varphi)}_{\text{surf}} \approx \varphi\, r_{\text{valence}}/3$.
        \item A candidate expression for Boltzmann’s constant $k_B$ is derived from the same set of integers, the electron’s minimal pulsation, the dynamical fraction of the proton’s relational sea, and a single coherence–leakage parameter $\varepsilon$. The resulting formulation reproduces the empirical order of magnitude of $k_B$ within a controlled and explicitly quantifiable relative error.
        \item An effective gravitational coupling is obtained by treating $\varepsilon$ as a minimal, hierarchically amplified leakage of coherence from composite nucleonic structures into the surrounding relational network. In this framework, one recovers a value for the gravitational constant that is compatible with observations when expressed in terms of $\hbar$, $c$, $m_p$, $\alpha$, and $\varepsilon$.
    \end{itemize}

    None of these constructions purports to yield exact predictions at the level of current experimental precision. Rather, the central claim is that several a priori conceptually distinct constants—$\alpha$, $k_B$, and the gravitational coupling—are associated with the \emph{same} underlying discrete architecture and governed by \emph{one} leakage parameter, with no additional free continuous parameters beyond the standard empirically determined mass ratio. This structural rigidity renders the framework highly susceptible to empirical falsification: increasing the stringency of either the combinatorial constraints or the experimental data can, in principle, decisively exclude it.
\end{enumerate}

Beyond these specific results, the PDL programme furnishes a unified and internally consistent framework for addressing several outstanding problems in fundamental physics, while systematically formulating all concepts and structures in an intrinsically discrete and combinatorial language, rather than assuming a priori any underlying geometric continuum.

\begin{itemize}[label=-]
    \item Proper time is reinterpreted as the discrete counting of cycles of internal coherence, with relativistic time dilation arising as a modification of the coherence budget required to maintain internal dynamical regimes across distinct relational environments.
    \item Gravitation is conceptualized as a variation in coherence cost within an emergent metric structure constructed from compatibility relations among coexisting closures, rather than as a fundamental field defined a priori on a fixed manifold.
    \item Black holes are described as extremal constraint configurations whose event horizon naturally functions as an information–bearing boundary within a relational network, thereby unifying Bekenstein–Hawking entropy with the internal surface structure already postulated for the proton.
    \item A speculative cosmological framework is proposed in which Big–Bang–like phenomena and extreme gravitational singularities are interpreted as global constraint events in a pre‑physical logical field, within which stable closures are both generated and recycled.
\end{itemize}

PDL does \emph{not} constitute a complete or all-encompassing formal system. Rather, it should be regarded as a rigorously specified framework that:
\begin{itemize}[label=-]
    \item is grounded in explicit axioms formulated on signed graphs rather than on continuum field theories;
    \item yields well-defined discrete structural models for the electron, proton, and other particles, characterized by integer-valued quantities that can be empirically and theoretically scrutinized;
    \item relates several physical constants to these discrete architectures in ways that are numerically non-trivial, while remaining amenable to subsequent refinement; and
    \item formulates hypotheses with sufficient rigor and specificity to allow empirical falsification through future high-precision measurements or more sophisticated combinatorial analyses.
\end{itemize}

For readers working in the areas of quantum foundations, quantum information theory, or mathematical physics, PDL should therefore not be viewed merely as an additional ontological framework, but rather as a concrete proposal that systematically investigates the extent to which the hypothesis can be sustained that coherence—rather than matter or spacetime geometry—constitutes the fundamental invariant. On this view, the empirically accessible world is interpreted as the unique realization of a minimal logic of relations that is capable of maintaining its own persistence.

\end{document}
